% This document demonstrates use of the LaTeX ``figshare-meta'' style.
% The style helps parse and typeset Figshare metadata.
%
% Passing a language parameter to the ``article'' document class is recommended.
% The language parameter will be inherited by the ``babel'' package used in ``figshare-meta''.
% Passing a language parameter typically improves typesetting elements such as line breaking.
% Here, typesetting is being enhanced for ``english'' content.
\documentclass[english]{article}
\usepackage{figshare-meta}   % For metadata extraction from Figshare.
\usepackage{parskip}   % For keeping a blank line between paragraphs with the `article' class.

\begin{document}

% Import and parse a target Figshare metadata file.
% JSON-formatted metadata is assumed.
% To typeset metadata for a different Figshare item, modify the file name, below, accordingly.
\FigshareMetaLoadJSON{../data/figshare-metadata-example.json}

\section*{\FigshareMetaGetTitle}

\textbf{Authors:} \FigshareMetaGetAuthorsFullName%

\FigshareMetaBeginHtmlBreaks%
\textbf{Description:} \FigshareMetaGetDescription%
\FigshareMetaEndHtmlBreaks%
\par

\textbf{Categories:} \FigshareMetaGetCategories%

\textbf{DOI:} \FigshareMetaGetDOI*%

\textbf{Figshare ID:} \FigshareMetaGetID%

\FigshareMetaTabulateFiles%

\end{document}
